\section{Intersection of local boundary spaces along a strip}%
\label{sec:peps_g_injective_strip_intersection}

\begin{lemma}[Averaged left inverse of an invariant MPS tensor]%
    \label{lem:peps_mps_g_injective_averaged_left_inverse}
    \lean{TNLean.PEPS.IsGInjective.exists_mpsLeftInverse}
    \leanok
    \uses{def:peps_mps_g_injective}
    If $A$ is $G$-injective for a representation $\rho$ of a finite
    group on a finite-dimensional virtual space, then
    $\rho(g)A^i\rho(g)^{-1}=A^i$ and there is a linear map $L_A$
    from physical vectors to virtual operators such that
    \[
        L_A\mathcal P(A)=\sigma,\qquad
        \sigma(X)=|G|^{-1}\sum_g\rho(g)X\rho(g)^{-1}.
    \]
    This is \cite[Definition~4.2 and Lemma~4.7]{Schuch2010PEPS}
    with the left inverse extended by group averaging.
\end{lemma}

\begin{proof}\leanok
    \uses{thm:peps_g_injective_left_inverse,
        thm:peps_mps_g_injective_source_form}
    Apply the averaged left-inverse characterization of $G$-injectivity.
    Invariance of the trace-pairing map is equivalent to conjugation
    invariance of its tensor coefficients.
\end{proof}

\begin{definition}[Lifted strip boundary spaces]%
    \label{def:peps_g_injective_strip_spaces}
    \lean{TNLean.PEPS.stripLeftSlice,
        TNLean.PEPS.stripRightSlice,
        TNLean.PEPS.stripLeftGroundSpace,
        TNLean.PEPS.stripRightGroundSpace,
        TNLean.PEPS.mem_stripLeftGroundSpace_iff,
        TNLean.PEPS.mem_stripRightGroundSpace_iff}
    \leanok
    Let $A^i,B^j,C^k$ be operator-valued tensors on a common
    finite-dimensional virtual space $\mathcal V$, with independent
    physical alphabets. Write
    \[
        \mathcal G_{AB}
        =\{(\operatorname{tr}[A^iB^jX])_{ij}:X\in\operatorname{End}\mathcal V\},
        \qquad
        \mathcal G_{BC}
        =\{(\operatorname{tr}[B^jC^kX])_{jk}:X\in\operatorname{End}\mathcal V\}.
    \]
    Define $\mathcal L$ to consist of three-block vectors whose every
    fixed-$k$ slice belongs to $\mathcal G_{AB}$, and $\mathcal R$
    to consist of those whose every fixed-$i$ slice belongs to
    $\mathcal G_{BC}$. Thus
    \[
        \mathcal L=\mathcal G_{AB}\otimes\mathcal P_C,
        \qquad
        \mathcal R=\mathcal P_A\otimes\mathcal G_{BC}.
    \]
    These are the local boundary spaces with the unused physical
    factor unrestricted, as in \cite[Theorem~4.8]{Schuch2010PEPS}.
\end{definition}

\begin{theorem}[Inhomogeneous strip intersection]%
    \label{thm:peps_g_injective_strip_intersection}
    \lean{TNLean.PEPS.IsGInjective.strip_intersection_property_of_invariant_left,
        TNLean.PEPS.IsGInjective.strip_groundSpace_intersection_of_invariant_left,
        TNLean.PEPS.IsGInjective.strip_groundSpace_intersection}
    \leanok
    \uses{def:peps_g_injective_strip_spaces,def:peps_mps_g_injective}
    Suppose a finite group acts on $\mathcal V$ through $\rho$, the
    middle and right blocks $B,C$ are $G$-injective, and the left block
    satisfies $\rho(g)A^i\rho(g)^{-1}=A^i$. Then
    \[
        \mathcal L\cap\mathcal R
        =\{(\operatorname{tr}[A^iB^jC^kX])_{ijk}
        :X\in\operatorname{End}\mathcal V\}.
    \]
    In particular this holds when all three blocks are $G$-injective.
    The endpoints may have different tensors and physical alphabets.
    This is the matrix-valued strip argument used in
    \cite[Theorems~4.8 and~5.4]{Schuch2010PEPS}; identifying an arbitrary
    graph's regional open contractions with these strip maps is a
    separate geometric step.
\end{theorem}

\begin{proof}\leanok
    \uses{lem:peps_mps_g_injective_left_inverse_traces,
        lem:peps_mps_g_injective_averaged_left_inverse,
        lem:peps_mps_g_injective_twirl,
        thm:peps_mps_g_injective_concatenation}
    A vector in the intersection has representations
    $\psi_{ijk}=\operatorname{tr}[A^iB^jM^k]
        =\operatorname{tr}[N^iB^jC^k]$.
    The two boundary tensors $M,N$ and the common boundary $X$ give
    % Formula: tr(A^i B^j M^k)=tr(N^i B^j C^k)=tr(A^i B^j C^k X).
    % Source: Papers/1001.3807/paper_v3.tex:1084--1120,
    % figs2/trAAM, figs2/trNAA, figs2/trAAA and prove-N-eq-BX-sym;
    % these are the inhomogeneous versions used by this theorem.
    % Ink-to-index: the three north legs are i,j,k in this order; all
    % horizontal joins and the returning arc are summed virtual indices.
    % M has the third physical leg, N the first, and X has no physical leg.
    % Boundary signature: each panel has (V,P)=(0,3).
    \begin{tenkzequation}
        \tenkzkernel
        \begin{tenkz}[rows={wire}, cols=3, bonds=none]
            \tn[name=stripA, ports={180:virtual,0:virtual,90:physical:$i$}]{A}
            \tn[at=(1,2), name=stripB,
                ports={180:virtual,0:virtual,90:physical:$j$}]{B}
            \tn[at=(1,3), name=stripM,
                ports={180:virtual,0:virtual,90:physical:$k$}]{M}
            \tnwire{stripA.0}{stripB.180}
            \tnwire{stripB.0}{stripM.180}
            \tnwire[route=arc]{stripM.0}{stripA.180}
        \end{tenkz}
        $=$
        \begin{tenkz}[rows={wire}, cols=3, bonds=none]
            \tn[name=stripN, ports={180:virtual,0:virtual,90:physical:$i$}]{N}
            \tn[at=(1,2), name=stripBtwo,
                ports={180:virtual,0:virtual,90:physical:$j$}]{B}
            \tn[at=(1,3), name=stripC,
                ports={180:virtual,0:virtual,90:physical:$k$}]{C}
            \tnwire{stripN.0}{stripBtwo.180}
            \tnwire{stripBtwo.0}{stripC.180}
            \tnwire[route=arc]{stripC.0}{stripN.180}
        \end{tenkz}
        $=$
        \begin{tenkz}[rows={wire}, cols=4, bonds=none]
            \tn[name=stripAthree,
                ports={180:virtual,0:virtual,90:physical:$i$}]{A}
            \tn[at=(1,2), name=stripBthree,
                ports={180:virtual,0:virtual,90:physical:$j$}]{B}
            \tn[at=(1,3), name=stripCthree,
                ports={180:virtual,0:virtual,90:physical:$k$}]{C}
            \tn[at=(1,4), name=stripX, skin=ring,
                ports={180:virtual,0:virtual}]{X}
            \tnwire{stripAthree.0}{stripBthree.180}
            \tnwire{stripBthree.0}{stripCthree.180}
            \tnwire{stripCthree.0}{stripX.180}
            \tnwire[route=arc]{stripX.0}{stripAthree.180}
        \end{tenkz}
    \end{tenkzequation}
    Apply the left inverse of the middle-right block. If $L_C$ is
    the left inverse of $C$ and $\sigma$ is group averaging on virtual
    operators, this gives
    \[
        \sigma(N^i)=XA^i,
        \qquad
        X=|G|\sum_k L_C|k\rangle\,\Delta\,\sigma(M^k).
    \]
    The invariant middle-right tensor permits replacing $N^i$ by
    $\sigma(N^i)$ in its trace pairing. Cyclicity of the trace then
    gives the three-block representation. Conversely, a three-block
    boundary $X$ gives $M^k=C^kX$ and $N^i=XA^i$.
\end{proof}

\begin{corollary}[Unique invariant boundary and intersection dimension]%
    \label{cor:peps_g_injective_strip_boundary}
    \lean{TNLean.PEPS.IsGInjective.existsUnique_invariant_boundary_of_mem_strip_intersection,
        TNLean.PEPS.IsGInjective.finrank_strip_groundSpace_intersection}
    \leanok
    \uses{thm:peps_g_injective_strip_intersection}
    If all three blocks are $G$-injective and all physical alphabets
    are finite, every vector in $\mathcal L\cap\mathcal R$ has a unique
    boundary operator $X$ commuting with $\rho(G)$. Consequently
    \[
        \dim(\mathcal L\cap\mathcal R)
        =\dim\operatorname{End}(\mathcal V)^G.
    \]
\end{corollary}

\begin{proof}\leanok
    \uses{thm:peps_mps_g_injective_concatenation,
        thm:peps_g_injective_invariants_range}
    Concatenation makes the three-block map $G$-injective. It therefore
    identifies the invariant virtual operator space with its physical
    range. Apply the intersection theorem.
\end{proof}
